% Reviewed source SHA256 (UTF-8/LF): 9fe78337c71646b3e9de33810af4b2985d24023aeac00e3f4c8008700c5ad5a2
% Options for packages loaded elsewhere
\PassOptionsToPackage{unicode}{hyperref}
\PassOptionsToPackage{hyphens}{url}
%
\documentclass[
  11pt,
]{article}
\usepackage{amsmath,amssymb}
\usepackage{setspace}
\usepackage{iftex}
\ifPDFTeX
  \usepackage[T1]{fontenc}
  \usepackage[utf8]{inputenc}
  \usepackage{textcomp} % provide euro and other symbols
\else % if luatex or xetex
  \usepackage{unicode-math} % this also loads fontspec
  \defaultfontfeatures{Scale=MatchLowercase}
  \defaultfontfeatures[\rmfamily]{Ligatures=TeX,Scale=1}
\fi
\usepackage{lmodern}
\ifPDFTeX\else
  % xetex/luatex font selection
\fi
% Use upquote if available, for straight quotes in verbatim environments
\IfFileExists{upquote.sty}{\usepackage{upquote}}{}
\IfFileExists{microtype.sty}{% use microtype if available
  \usepackage[]{microtype}
  \UseMicrotypeSet[protrusion]{basicmath} % disable protrusion for tt fonts
}{}
\makeatletter
\@ifundefined{KOMAClassName}{% if non-KOMA class
  \IfFileExists{parskip.sty}{%
    \usepackage{parskip}
  }{% else
    \setlength{\parindent}{0pt}
    \setlength{\parskip}{6pt plus 2pt minus 1pt}}
}{% if KOMA class
  \KOMAoptions{parskip=half}}
\makeatother
\usepackage{xcolor}
\usepackage[margin=1in]{geometry}
\setlength{\emergencystretch}{3em} % prevent overfull lines
\providecommand{\tightlist}{%
  \setlength{\itemsep}{0pt}\setlength{\parskip}{0pt}}
\setcounter{secnumdepth}{-\maxdimen} % remove section numbering
\usepackage{amsmath,amssymb}
\usepackage{microtype}
\usepackage{xurl}
\allowdisplaybreaks[1]
\setlength{\emergencystretch}{2em}
\AtBeginDocument{\hypersetup{colorlinks=true,linkcolor=black,urlcolor=black,citecolor=black}}

\widowpenalty=10000

\clubpenalty=10000

\brokenpenalty=10000
\ifLuaTeX
  \usepackage{selnolig}  % disable illegal ligatures
\fi
\usepackage{bookmark}
\IfFileExists{xurl.sty}{\usepackage{xurl}}{} % add URL line breaks if available
\urlstyle{same}
\hypersetup{
  pdftitle={Conditioning of random cyclic-shift Krylov compressions},
  hidelinks,
  pdfcreator={LaTeX via pandoc}}

\title{Conditioning of random cyclic-shift Krylov compressions}
\author{Matthew J. Colbrook\\\small Department of Applied Mathematics and Theoretical Physics\\\small University of Cambridge, Cambridge, United Kingdom\\\small\href{mailto:m.colbrook@damtp.cam.ac.uk}{m.colbrook@damtp.cam.ac.uk}}
\date{11 September 2026}

\usepackage{xurl}
\setlength{\emergencystretch}{3em}
\hypersetup{pdfauthor={Matthew J. Colbrook}}
\begin{document}
\maketitle

\noindent\textbf{Independently reviewed resolution.} Authorship is recorded at the submitter's request. The supplied AI-generated proof and computational support are disclosed in the submission record. The exact canonical target is resolved.
The dated \href{https://github.com/MColbrook/OpenProblemsInNLA/blob/codex/colbrook-round3-submissions/references/colbrook-round3-2026-09-11/verification/reviews/IE-10-review.md}{independent agent review} records the subsequent checks and supersedes original draft remarks about pending review. This is not external human peer review or formal certification.
\par\medskip
% BEGIN REVIEWED BODY


\setstretch{1.04}
\subsection{1. Claim and scope}\label{claim-and-scope}

This manuscript proposes an affirmative resolution of repository problem
\textbf{IE-10}, with exactly its complex-sphere starting-vector
model.\footnote{Source 1, the displayed probability target and
  starting-vector model.} The original workshop question specifically
mentions the cyclic shift; a counterexample for general nonnormal inputs
does not address this case.\footnote{Source 2, Section 3.3, Problem 3.5
  and its 20 August 2026 update.}

\textbf{Theorem 1 (cyclic shift).} Let \(n\geq3\), let \(2\leq k<n\),
and let \(C_n\) be the cyclic shift on \(\mathbb C^n\). Draw \(b\)
uniformly from the complex unit sphere. Let \(Q\) have orthonormal
columns spanning \[
\operatorname{span}\{b,C_nb,\ldots,C_n^{k-1}b\},
\qquad H_k=Q^*C_nQ.
\] For a diagonalizable matrix define \[
\kappa_V(H)=\inf_{H=VDV^{-1},\ D\text{ diagonal}}\|V\|_2\|V^{-1}\|_2,
\] and set this quantity to infinity otherwise. Then \[
\mathbb E\,\kappa_V(H_k)\leq17n^2k. \tag{1}
\] Consequently, for every \(0<\eta<1\), \[
\Pr\left\{\kappa_V(H_k)\leq\frac{17}{\eta}n^2k\right\}\geq1-\eta. \tag{2}
\] In particular, \textbf{\(C=1700\) and \(c=3\)} satisfy the
repository's uniform \(0.99\) target.

Two probability arguments are given. Section 5 already proves the
repository's existence claim, with the weaker constants
\(C=2{,}000{,}000\) and \(c=5\), using discriminant avoidance and
Cauchy's estimate. Sections 6--7 improve this to (1)--(2) by controlling
the total variation of a rank-one root locus. Thus the affirmative
answer does not rely solely on the sharper root-locus estimate.

No assertion is made for real-sphere starts, arbitrary nonnormal input
matrices, or finite-precision construction of an eigenbasis.

\subsection{2. Exact weighted-polynomial
model}\label{exact-weighted-polynomial-model}

We work first with arbitrary distinct points \(\zeta_1,\ldots,\zeta_n\)
on the unit circle. Write \[
d=\min_{i\ne j}|\zeta_i-\zeta_j|>0,
\] and let \(w_1,\ldots,w_n\) be independent exponential random
variables of mean one. On the sample space
\(\mathbb C^{\{\zeta_1,\ldots,\zeta_n\}}\), use \[
\langle f,g\rangle_w=\sum_{j=1}^n w_j\overline{f(\zeta_j)}g(\zeta_j).
\] Let \(\mathcal P_{k-1}\) be the evaluated polynomials of degree less
than \(k\), let \(P_w\) be its orthogonal projection, and set \[
H_w=P_wM_z\big|_{\mathcal P_{k-1}},\qquad (M_zf)(\zeta_j)=\zeta_jf(\zeta_j).
\] Positive weights and distinct nodes make the evaluation map injective
on this polynomial space. Since \(M_z\) is unitary, \(H_w\) is a
contraction.

The general estimate proved below is \[
\mathbb E\,\kappa_V(H_w)\leq2k+\frac{64nk}{d}. \tag{3}
\] Here and below an operator condition number is measured in its stated
Hilbert-space norm, equivalently in any orthonormal coordinates.

To obtain the cyclic model, diagonalize \(C_n\) by a unitary Fourier
matrix. A uniformly distributed complex unit vector can be represented
as \(g/\|g\|_2\), where the coordinates of \(g\) are independent
standard complex Gaussians. This follows directly from the unitary
invariance of the density proportional to \(\exp(-\|g\|_2^2)\). In polar
coordinates, \(|g_j|^2\) are independent mean-one exponentials. The
Krylov vectors in spectral coordinates are the evaluations of
\(1,z,\ldots,z^{k-1}\) multiplied coordinatewise by \(g_j/\|g\|_2\).
Their phases are removed by a diagonal unitary commuting with the
spectral matrix. The common normalization of the weights does not change
an orthogonal compression. Thus \(H_k\) is unitarily equivalent to
\(H_w\) with the \(n\)th roots of unity as nodes. This also proves the
almost-sure dimension assertion.

For coefficient calculations put \[
a_j=(1,\zeta_j,\ldots,\zeta_j^{k-1}),\qquad
G=\sum_jw_ja_j^*a_j,\qquad F=\sum_jw_j\zeta_ja_j^*a_j.
\] The coefficient representation of \(H_w\) is \(G^{-1}F\), with
coefficient inner product \(x^*Gy\).

\subsection{3. Conditioning is controlled by weight
sensitivity}\label{conditioning-is-controlled-by-weight-sensitivity}

\subsubsection{3.1 A conjugation that matches left and right coordinate
magnitudes}\label{a-conjugation-that-matches-left-and-right-coordinate-magnitudes}

Define an antilinear map on the full sample space by \[
(Jf)(\zeta_j)=\zeta_j^{k-1}\overline{f(\zeta_j)}.
\] It is an antiunitary involution. For a polynomial of degree less than
\(k\), it reverses and conjugates the coefficients, so it preserves
\(\mathcal P_{k-1}\). Hence \(J\) commutes with \(P_w\). Directly, \[
JM_zJ=M_z^*,\qquad JH_wJ=H_w^*. \tag{4}
\] Suppose \(\lambda_i\) is simple and \(r_i\) is a unit right
eigenvector of \(H_w\). Then \(\ell_i=Jr_i\) is a unit left eigenvector,
meaning \(H_w^*\ell_i=\overline{\lambda_i}\ell_i\). Put \[
\alpha_i=\langle\ell_i,r_i\rangle_w,\qquad
\kappa_i=\frac{1}{|\alpha_i|},\qquad
p_{ij}=w_j|r_i(\zeta_j)|^2.
\] Simplicity gives \(\alpha_i\ne0\). The number \(\kappa_i\) is the
norm of the rank-one spectral projector. Also \[
\sum_jp_{ij}=1,\qquad
w_j|\ell_i(\zeta_j)|^2=p_{ij}. \tag{5}
\]

\subsubsection{3.2 An exact derivative
identity}\label{an-exact-derivative-identity}

Differentiate the generalized eigenvalue equation
\(Fr_i=\lambda_iGr_i\). Multiplication by the generalized left
eigenvector cancels the derivative of \(r_i\), giving \[
\frac{\partial\lambda_i}{\partial w_j}
=\frac{(\zeta_j-\lambda_i)\overline{\ell_i(\zeta_j)}r_i(\zeta_j)}{\alpha_i}.
\] Therefore \[
\sum_jw_j\left|\frac{\partial\lambda_i}{\partial w_j}\right|
=\kappa_i\sum_jp_{ij}|\zeta_j-\lambda_i|. \tag{6}
\] This identity is taken at simple-spectrum weight vectors; those
vectors have full measure by Section 4.

\textbf{Lemma 2 (deterministic sensitivity bound).} At every positive
simple-spectrum weight vector, \[
\kappa_i\leq
\max\left\{2,\frac8d\sum_jw_j\left|\frac{\partial\lambda_i}{\partial w_j}\right|\right\}
\leq2+\frac8d\sum_jw_j\left|\frac{\partial\lambda_i}{\partial w_j}\right|. \tag{7}
\]

\emph{Proof.} Each summand in
\(\alpha_i=\sum_jw_j\overline{\ell_i(\zeta_j)}r_i(\zeta_j)\) has modulus
\(p_{ij}\). The reverse triangle inequality implies \[
|\alpha_i|\geq2\max_jp_{ij}-1.
\] If \(\kappa_i>2\), then \(\max_jp_{ij}<3/4\). At most one node can
lie at distance strictly less than \(d/2\) from \(\lambda_i\).
Consequently at least one quarter of the probability mass in (5) is at
distance at least \(d/2\), and \[
\sum_jp_{ij}|\zeta_j-\lambda_i|\geq d/8.
\] Now use (6). The case \(\kappa_i\leq2\) is immediate. \(\square\)

\subsubsection{3.3 From individual projectors to an
eigenbasis}\label{from-individual-projectors-to-an-eigenbasis}

\textbf{Lemma 3 (balanced column scaling).} For any matrix with simple
spectrum, \[
\kappa_V(H)\leq\sum_{i=1}^k\kappa_i, \tag{8}
\] where the \(\kappa_i\) are the norms of its rank-one spectral
projectors.

\emph{Proof.} Use orthonormal ambient coordinates. Let \(R\) have unit
right eigenvectors as columns. The \(i\)th row of \(R^{-1}\) has norm
\(\kappa_i\): it is the dual left eigenvector with pairing one. Define
\(D=\operatorname{diag}(\sqrt{\kappa_1},\ldots,\sqrt{\kappa_k})\) and
\(V=RD\). Then \[
\|V\|_F^2=\sum_i\kappa_i,
\qquad
\|V^{-1}\|_F^2=\sum_i\frac{\kappa_i^2}{\kappa_i}=\sum_i\kappa_i.
\] The operator norms are bounded by the Frobenius norms, and \(V\) is
still an eigenbasis. \(\square\)

\subsection{4. Simplicity and rank-one conditional
pencils}\label{simplicity-and-rank-one-conditional-pencils}

The discriminant in \(z\) of \[
p(z,w)=\det(zG(w)-F(w))
\] is a polynomial in the weights and is not identically zero. To see
this, allow boundary weights supported positively on exactly \(k\)
distinct nodes. The resulting \(G\) remains positive definite. The
square evaluation matrix on those nodes is invertible, and \(G^{-1}F\)
is similar to the diagonal matrix of those \(k\) distinct nodes. Its
discriminant is nonzero. Thus the discriminant polynomial is nonzero,
and its zero set in real weight space has measure zero.

Fix an index \(j\) and condition on the other weights. For almost every
such conditioning the discriminant, considered as a polynomial in
\(t=w_j\), is not identically zero: at least one of its coefficient
polynomials in the remaining weights is nonzero.

Write \(a=a_j\) and \(\zeta=\zeta_j\). The conditional pencil is \[
G(t)=G_0+ta^*a,\qquad F(t)=F_0+t\zeta a^*a.
\] Because \(k<n\), the other \(n-1\) positive-weight nodes span the
polynomial space, so \(G_0\succ0\). By the rank-one determinant
identity, or multilinearity of the determinant, \[
p(z,t)=\det(zG(t)-F(t))=p_0(z)+tp_1(z),\qquad
\deg p_0,\deg p_1\leq k. \tag{9}
\] Its leading coefficient in \(z\) is \(\det G(t)>0\) for real \(t>0\).
Its discriminant \(\Delta(t)\) has degree at most \(2k-2\), since the
discriminant is homogeneous of that degree in the coefficients of a
degree-\(k\) polynomial and each coefficient in (9) is affine in \(t\).

For positive real \(t\), every root lies in the closed unit disk. Except
at the finitely many positive zeros of \(\Delta\), the roots admit
locally analytic labels. All sums over the roots used below are
independent of the chosen labeling.

\subsection{5. A direct probability proof using discriminant
avoidance}\label{a-direct-probability-proof-using-discriminant-avoidance}

This section independently proves a polynomial bound sufficient for
IE-10, without the length estimate of Section 6.

\textbf{Proposition 4 (weaker uniform tail bound).} In the cyclic model,
for \(0<\eta<1\), \[
\Pr\left\{\kappa_V(H_k)\leq\frac{192}{\eta^2}n^3k^2\right\}\geq1-\eta. \tag{10}
\] In particular, \(C=2{,}000{,}000\), \(c=5\) work at success
probability \(0.99\).

\emph{Proof.} Put \(\varepsilon=\eta/(16nk)\). For one conditional
pencil, call a positive real \(t_0\) bad if \(t_0\leq2\varepsilon\) or
if it is within distance \(2\varepsilon\) of a complex zero of
\(\Delta\). There are at most \(2k-2\) such zeros, counted with
multiplicity. A disk of radius \(2\varepsilon\) cuts the real line in an
interval of length at most \(4\varepsilon\). Since the exponential
density is bounded by one, the conditional bad probability is at most \[
2\varepsilon+4(2k-2)\varepsilon\leq8k\varepsilon.
\] The exceptional conditionings from Section 4 have probability zero.
Union over the \(n\) coordinates therefore costs at most
\(8nk\varepsilon=\eta/2\).

At a good \(t_0\), every root is holomorphic on \(|t-t_0|<\varepsilon\).
Indeed, the discriminant does not vanish there, and \(G(t)\) is
invertible in \(\Re t>0\). For an eigenpair of the complex conditional
pencil, set \[
g=x^*G_0x>0,\quad m=x^*F_0x,\quad s=|ax|^2\geq0.
\] The positive-weight representation of \(F_0\) gives \(|m|\leq g\).
Multiplying the generalized eigenvalue equation by \(x^*\) gives \[
z=\frac{m+t\zeta s}{g+ts}.
\] Throughout the disk, \(\Re t>t_0/2\) and \(|t|<3t_0/2\), so \[
|z|\leq\frac{g+(3t_0/2)s}{g+(t_0/2)s}\leq3.
\] Cauchy's estimate yields, for every eigenvalue and every good
coordinate, \[
\left|\frac{\partial\lambda_i}{\partial w_j}\right|\leq3/\varepsilon. \tag{11}
\]

Let \(W=\sum_jw_j\). Since \(\mathbb EW=n\), Markov's inequality gives
\[
\Pr\{W>2n/\eta\}\leq\eta/2.
\] On the joint good event, of probability at least \(1-\eta\),
equations (7), (8), and (11) give \[
\kappa_V(H_k)
\leq k\max\left\{2,\frac{24W}{d\varepsilon}\right\}
\leq\frac{768n^2k^2}{d\eta^2}.
\] For cyclic nodes, \(d\geq4/n\); the bound dominates \(2k\) and proves
(10). \(\square\)

\subsection{6. Total variation of the conditional root
locus}\label{total-variation-of-the-conditional-root-locus}

\textbf{Lemma 5 (bounded root-locus length).} For almost every
conditioning in Section 4, \[
\sum_{i=1}^k\int_0^\infty|\lambda_i'(t)|\,dt\leq8k, \tag{12}
\] where the integrals are taken over the simple-spectrum intervals.
Values at the finite collision set are irrelevant.

\emph{Proof.} If \(p_1\) vanishes identically, all roots are constant.
Otherwise define the real polynomial \[
q(x,y)=\Im\left(p_0(x+iy)\overline{p_1(x+iy)}\right),
\qquad \deg q\leq2k.
\] A moving root, away from a common root of \(p_0\) and \(p_1\),
satisfies \[
t=-\frac{p_0(z)}{p_1(z)}\in(0,\infty),\qquad q(\Re z,\Im z)=0. \tag{13}
\] Common roots give constant branches. A moving branch can meet one
only at an exceptional collision parameter, so these points do not
affect the integrals.

If \(q\) vanishes identically, the holomorphic rational function
\(p_0/p_1\) is real valued on its domain and is therefore constant.
Equation (9), together with the nonvanishing leading coefficient for
positive \(t\), again implies that all roots are constant. We may
therefore assume \(q\not\equiv0\).

For all but finitely many real \(u\), the polynomial \(q(u,y)\) is
nonzero and has at most \(2k\) real roots. Each moving-root point
\(u+iy\) with \(p_1(u+iy)\ne0\) determines at most one parameter \(t\)
through (13), and at a simple-spectrum parameter it belongs to only one
root. Consequently, for almost every \(u\), the number of moving-root
crossings of the vertical line \(\Re z=u\), counted with their parameter
preimages, is at most \(2k\). The same statement holds for horizontal
lines.

For clarity, this crossing count bounds variation, not merely the
set-theoretic image of a curve. On each simple-spectrum interval a root
is analytic. Split each real or imaginary coordinate into its
monotonicity intervals, ignoring constant coordinates, and apply
one-dimensional change of variables on those intervals. Countably many
intervals suffice; nonnegative integrands allow summation by Tonelli's
theorem. Since all roots are in the closed unit disk, their real and
imaginary coordinates are in \([-1,1]\). Thus \[
\sum_i\int_0^\infty|\Re\lambda_i'(t)|\,dt
\leq\int_{-1}^1 2k\,du=4k,
\] and likewise \[
\sum_i\int_0^\infty|\Im\lambda_i'(t)|\,dt\leq4k.
\] The exceptional coordinate values corresponding to collisions,
constant branches, or lines contained in \(q=0\) form a null set.
Finally, \(|z'|\leq|\Re z'|+|\Im z'|\) gives (12). \(\square\)

\subsection{7. Proof of the expectation
bound}\label{proof-of-the-expectation-bound}

Condition on all weights except \(w_j=t\). Its remaining density is
\(e^{-t}\) on \((0,\infty)\). Lemma 5 gives \[
\begin{aligned}
\mathbb E\left[\sum_iw_j\left|\frac{\partial\lambda_i}{\partial w_j}\right|
\,\middle|\,(w_\ell)_{\ell\ne j}\right]
&=\int_0^\infty t e^{-t}\sum_i|\lambda_i'(t)|\,dt\\
&\leq8k,
\end{aligned} \tag{14}
\] where we used only \(te^{-t}\leq1\). Exceptional conditionings have
probability zero. Sum (14) over \(j\), and then use (7)--(8) and
Tonelli's theorem: \[
\mathbb E\,\kappa_V(H_w)
\leq2k+\frac8d\mathbb E\sum_{i,j}w_j
\left|\frac{\partial\lambda_i}{\partial w_j}\right|
\leq2k+\frac{64nk}{d}.
\] This proves (3). For the cyclic shift, \[
d=2\sin(\pi/n)\geq\frac4n,
\] using concavity of sine on \([0,\pi/2]\). Therefore \[
\mathbb E\,\kappa_V(H_k)\leq2k+16n^2k\leq17n^2k.
\] Markov's inequality proves Theorem 1.

\textbf{Corollary 6 (one starting vector, all dimensions).} With the
same random \(b\) used for every \(k\), no independence between the
compressions is needed to conclude \[
\Pr\left\{\kappa_V(H_k)\leq\frac{17}{2\eta}n^4
\text{ for every }2\leq k<n\right\}\geq1-\eta.
\] Indeed, the expectation of the sum of their condition numbers is at
most \(17n^2\sum_{k=2}^{n-1}k\leq(17/2)n^4\).

\subsection{8. Checks, limitations, and
provenance}\label{checks-limitations-and-provenance}

The accompanying code checks the weighted model against an orthonormal
Krylov construction, the matched left/right magnitudes, the eigenvalue
derivative identity, balanced column scaling, finite-difference
derivatives, and sampled conditional root paths. A separate
140-decimal-digit diagnostic includes weight ratios of \(10^{24}\),
closely clustered nodes, and cyclic weights within \(10^{-35}\) of the
defective uniform example. These stress tests are numerical, not
interval certificates. An exact \(n=3\), \(k=2\) example with weights
\((t,1,1)\) gives \[
G=\begin{pmatrix}t+2&t-1\\t-1&t+2\end{pmatrix},\qquad
F=\begin{pmatrix}t-1&t-1\\t+2&t-1\end{pmatrix},
\] \[
\det(zG-F)=3\big((2t+1)z^2+(1-t)z+(1-t)\big),
\] with discriminant \(81(t-1)(t+1/3)\). This illustrates why an
almost-sure argument must allow exceptional defective weights; the point
\(t=1\) is not silently excluded by an unproved uniform deterministic
bound.

This example also permits an exact geometric check of Lemma 5. For
\(0<t<1\), the roots lie on the circle \(|z+1|=1\) and each moves
monotonically from angle \(\pm\pi/3\) to angle zero about its center.
For \(t>1\), one root moves from zero to \(1\) and the other from zero
to \(-1/2\). Indeed, \[
\frac{dt}{dz}=\frac{3z(z+2)}{(2z^2-z-1)^2},
\] which has the required sign on each real interval. The total
root-path length is therefore \(2\pi/3+3/2<16=8k\). This is a check on
one explicit pencil, not a substitute for the general crossing-count
proof.

The tests are diagnostics and exact algebraic checks, not an independent
proof of the probability theorem. The proof uses randomness only in the
starting vector through independent exponential spectral weights. It
never substitutes an independent Haar subspace for the Krylov subspace.
Both proofs retain the restrictions \(k<n\) and complex-sphere
randomness.

This is an AI-generated proposed argument, prepared for mathematical
review on 11 September 2026. It has not been independently refereed or
formally verified. No assertion of priority is made. The proposed
repository status is \textbf{Solution claimed}, not independently
established \textbf{Solved}.

\subsection{Sources}\label{sources}

\begin{enumerate}
\def\labelenumi{\arabic{enumi}.}
\item
  OpenProblemsInNLA.
  \href{https://github.com/ajt60gaibb/OpenProblemsInNLA/blob/main/eigenvalues-and-inverse-problems/IE-10/README.md}{IE-10:
  Conditioning of a random Krylov compression of a cyclic shift}.
  Problem statement and audit update of 10 September 2026. Accessed 11
  September 2026.
\item
  Noah Amsel et
  al.~\href{https://arxiv.org/html/2602.05394v3}{\emph{Linear Systems
  and Eigenvalue Problems: Open Questions from a Simons Workshop}}.
  arXiv:2602.05394v3, 21 August 2026. Section 3.3, Problem 3.5 and its
  update. This source identifies the question and scope; it is not the
  source of the proposed proof.
\end{enumerate}

\end{document}

% END REVIEWED BODY
